\section{Exact coefficient two for separated prime bands}
\label{sec:separated-bands}

The following result permits a growing, possibly dense small-prime support.
Its source restriction is that every source outside the pure small-prime
part contains exactly one large prime occurrence. The proof assigns targets
in disjoint exact smooth-part classes and then uses one global repair phase.
It does not combine separate covers of the two source parts.

Let $A\subseteq\{2,\ldots,M\}$ be a set of $n$ distinct integers and let
$I=\{x+1,\ldots,x+M\}$, where $x\in\mathbb Z_{\ge0}$. The length $M$
is an ambient bound and need not equal $\max A$. Let $2\le Y<Z$, let $Q$
be a subset of the primes at most $Y$, and let $R$ be a finite set of primes
at least $Z$. Write
\[
 Q^\Pi=\prod_{q\in Q}q,
 \qquad \vartheta(Q)=\prod_{q\in Q}(1-1/q),
 \qquad \pi_R(T)=\#\{r\in R:r\le T\}.
\]
Both empty products equal one, and one is regarded as $Q$-smooth. Partition
the sources as $A=A_Q\sqcup A_R$, where each source in $A_Q$ is $Q$-smooth
and each source in $A_R$ has the form
\begin{equation}
 a=d(a)r(a),\qquad d(a)\ \text{$Q$-smooth},\quad r(a)\in R.
 \label{eq:sep-source-form}
\end{equation}
Thus the large prime has exponent one; sources with two large prime
occurrences are not included in this hypothesis.

We use the decreasing error function $e(\kappa)\to0$ from
Lemma~\ref{an:sieve}. In particular, for every real interval $J$ of length
$T\ge Y^\kappa$, uniformly over its shift and the choice of $Q$,
\begin{equation}
 \#\{m\in J:(m,Q^\Pi)=1\}
 \ge (1-e(\kappa))\vartheta(Q)T,
 \qquad \kappa\ge4.
 \label{eq:sep-sieve}
\end{equation}
This is the fundamental-lemma input of \cite{Lagarias,GranvilleSound}; it
includes nonintegral interval lengths. We also use absolute constants
$c_0,C_0>0$ in the standard Mertens and Chebyshev bounds
\begin{equation}
 \prod_{p\le Y}(1-1/p)\ge\frac{c_0}{\log Y},
 \qquad \pi(T)\le\frac{C_0T}{\log T}
 \quad(Y,T\ge2).
 \label{eq:sep-prime-estimates}
\end{equation}
No numerical values for the resulting separation exponent are asserted.

\begin{proposition}[A finite support comparison]
\label{prop:sep-finite}
Suppose $\kappa\ge4$, $e(\kappa)<1$, $Z\ge Y^\kappa$, and
\begin{equation}
 (1-e(\kappa))\vartheta(Q)
 \ge \sup_{T\ge Z}\frac{\pi_R(T)}T.
 \label{eq:sep-finite-condition}
\end{equation}
Then, for every source set satisfying \eqref{eq:sep-source-form},
\begin{equation}
 \tau(A,I)\le 2|A_R|+\sum_{a\in A_Q}\Omega(a).
 \label{eq:sep-cover-bound}
\end{equation}
For finite $R=\{r_1<\cdots<r_s\}\ne\varnothing$, the supremum in
\eqref{eq:sep-finite-condition} is $\max_{1\le j\le s}j/r_j$; for empty
$R$ it is zero.
\end{proposition}

\begin{proof}
For each distinct $d$ occurring in \eqref{eq:sep-source-form}, define
\[
 m_d=\#\{a\in A_R:d(a)=d\},\qquad T_d=M/d,
 \qquad B_d=\{b\in I:d\mid b,\ (b/d,Q^\Pi)=1\}.
\]
Every $b\in B_d$ has exact $Q$-part $d$: the product of all its prime
powers at primes in $Q$ equals $d$. Hence the sets $B_d$ are pairwise
disjoint, including the case $d=1$. Distinct sources with this cofactor
have distinct primes $r\le T_d$, so
\begin{equation}
 m_d\le\pi_R(T_d),\qquad T_d\ge Z.
 \label{eq:sep-fibre}
\end{equation}
The multiples of $d$ in $I$, after division by $d$, are the integer points
of $(x/d,(x+M)/d]$, which has real length $T_d$. Therefore
\eqref{eq:sep-sieve} and \eqref{eq:sep-finite-condition} give
\[
 |B_d|\ge(1-e(\kappa))\vartheta(Q)T_d
 \ge\pi_R(T_d)\ge m_d.
\]
Choose $m_d$ distinct targets from each $B_d$ and assign them to the
corresponding sources. Since the bins are disjoint, all $|A_R|$ assigned
targets are distinct. Their product is divisible by
$\prod_{a\in A_R}d(a)$.

The sum of the remaining positive prime-valuation deficits relative to
$\prod A$ is at most
\[
 |A_R|+\sum_{a\in A_Q}\Omega(a).
\]
There is one possible unit for each removed prime $r(a)$, together with
all units of the pure $Q$ sources; extra factors already supplied by the
chosen targets only help. Since $\prod A\mid M!\mid\prod I$, each still
missing unit can be supplied by an unused target. Such a repair decreases
the total positive deficit by at least one. Beginning with $|A_R|$ assigned
targets, this one global repair phase proves \eqref{eq:sep-cover-bound}.

Finally, $\pi_R(T)/T$ decreases between its jumps at the primes of $R$,
which gives the displayed finite expression for its supremum.
\end{proof}

For a particular source set, the proof only needs
$(1-e(\kappa))\vartheta(Q)T_d\ge m_d$ for its actual fibres. Condition
\eqref{eq:sep-finite-condition} is a support-only sufficient comparison.

\begin{theorem}[Absolute power separation]
\label{thm:sep-power}
There is an absolute constant $\kappa_0\ge4$ such that whenever
$Y\ge2$ and $Z\ge Y^{\kappa_0}$, every source set of the form
\eqref{eq:sep-source-form} satisfies \eqref{eq:sep-cover-bound}, for every
ambient $M\ge\max A$ and every integer shift $x\ge0$. In particular, if
$A_Q=\varnothing$, then $\tau(A,I)\le2n$.
\end{theorem}

\begin{proof}
Choose an absolute $\kappa_0$ such that
\[
 e(\kappa_0)\le\tfrac12,
 \qquad \kappa_0\ge\frac{2C_0}{c_0}.
\]
This is possible because $e(\kappa)\to0$. For every $T\ge Z\ge
Y^{\kappa_0}$, the subset relation for $Q$ and
\eqref{eq:sep-prime-estimates} give
\[
 (1-e(\kappa_0))\vartheta(Q)
 \ge\frac{c_0}{2\log Y}
 \ge\frac{C_0}{\log T}
 \ge\frac{\pi_R(T)}T.
\]
Proposition~\ref{prop:sep-finite} applies.
\end{proof}

The constant $\kappa_0$ is not computed or claimed to be practical. The
theorem is a finite statement for all parameters meeting its separation
condition. It implies the proposed asymptotic class $Z\to\infty$,
$\log Y=o(\log Z)$; for example, $Y=M^{o(1)}$ and $Z=M^\alpha$ for any
fixed $\alpha>0$ eventually satisfy the condition. The restriction to one
large prime occurrence per source remains essential to the stated result.

\begin{corollary}[Pure small-prime sources]
\label{cor:sep-pure}
Under the hypotheses of Theorem~\ref{thm:sep-power},
\begin{equation}
 \tau(A,I)\le2n+\sum_{a\in A_Q}(\Omega(a)-2)_+.
 \label{eq:sep-excess}
\end{equation}
Thus $\tau(A,I)\le2n$ if every source in $A_Q$ has at most two prime
factors with multiplicity. More generally, it is enough that
$\sum_{a\in A_Q}\Omega(a)\le2|A_Q|$. Along any sequence with
$n\to\infty$ and $\sum_{a\in A_Q}(\Omega(a)-2)_+=o(n)$, the upper
bound is $(2+o(1))n$, uniformly in the shift.
\end{corollary}

\begin{proof}
Subtract $2n=2|A_R|+2|A_Q|$ from \eqref{eq:sep-cover-bound}, and bound
each possibly negative term $\Omega(a)-2$ by its positive part.
\end{proof}

This is an upper-bound result; no matching lower coefficient is asserted.
An unrestricted pure $Q$ component may have large excess in
\eqref{eq:sep-excess}, and sources with multiple $R$ occurrences remain
outside the theorem. Consequently the argument does not settle the
unrestricted exact or asymptotic problem.
